\section{Policy}\label{sec:policy}

\subsection{Economic Prosperity Deal}\label{sec:policy-epd}

The paired full-minus-EPD contrast is withdrawn as a quantitative result (Section~\ref{sec:results-epd}) and reported in the online supplement: each paired difference is comparable to or smaller than its assignment dispersion, and the underlying sector calibration fails its rank-order check, so no quantitative value should be placed on the deal figures. The design---identical downstream machinery under two declared tariff schedules---remains the template for pricing future schedule changes once a corroborated sector calibration exists.

\subsection{Pharmaceutical channel}\label{sec:policy-pharma}

The EPD scenario sets the pharmaceutical employment shock to zero, but this does not make the exemption free. Its VPAG and NHS pricing concessions sit outside the labour-income microsimulation. A complete policy appraisal would cost those health-budget and price effects alongside the direct tax--benefit scenario.

\subsection{Cushioning policy}\label{sec:policy-reforms}

The simulations support a limited conclusion: extensive-margin protection depends more on means-tested eligibility and claiming than the income-tax response to an intensive earnings loss. The full-draw take-up grid shows the claiming-rate assumption is itself a binding parameter---moving the new-entitlement rate from 0.55 to 1.00 shifts displacement cushioning by \TakeupConventionSpread\ percentage points---and the gates the model cannot exercise, the capital test and take-up among long-standing eligible non-claimants, would push measured cushioning further down. Automatic enrolment after redundancy, contribution-based unemployment insurance and wage insurance are therefore candidate counterfactuals for future analysis, not policy recommendations established by the present results. Their costs and incidence require dedicated simulations and financing assumptions.

The repository also supplies reusable interfaces for the next policy exercises.
Observable household and worker characteristics can tilt shock probabilities
while preserving each sector's aggregate earnings loss, and a separate
price-channel interface converts expenditure-weighted price changes into real
disposable income. Neither interface supplies causal coefficients: the former
requires externally estimated heterogeneity and the latter requires household
expenditure data. They are extensions to populate and validate, not additional
headline estimates in this paper.
